\scanpage{321}

\unsure{An} element of $\operatorname{Discr\acute{e}t}^{\unsure{\prime}}(\pi)$ is called a
class of discretifications of $\pi$\inserted{\unsure{or} \unsure{prediscretification} of $\pi$}.
If $\pi_0$ is a discretification, we denote its class by
$\pi_0^{\unsure{\natural}}$. The set of classes of discretifications on $\pi$
is a homogeneous \unsure{space} under $\widehat{\widehat{\unsure{\mathfrak{S}|\mathcal{G}}}}(\pi)$, the stabilizer
of $\pi_0^{\unsure{\natural}}$ \unsure{is identified with|being identified with} : the closure of
$\unsure{\mathfrak{S}|\mathcal{G}}(\pi_0)=\operatorname{Aut}_{\mathrm{lac}}(\pi_0)$ in
$\widehat{\widehat{\unsure{\mathfrak{S}|\mathcal{G}}}}(\pi)$, or again :
$\widehat{\unsure{\mathfrak{S}|\mathcal{G}}}(\pi_0)$.
The latter always contains $\pi$.
\[
\operatorname{Aut}_{\mathrm{lac}}(\pi,\pi_0^{\unsure{\natural}})
\simeq\widehat{\unsure{\mathfrak{S}|\mathcal{G}}}(\pi_0)
\]

\deleted{We have} for every discretification $\pi_0^{\unsure{\natural}}$, the subgroup of
$\widehat{\widehat{\unsure{\mathfrak{S}|\mathcal{G}}}}(\pi)$ of \unsure{automorphisms} with loops which fix
$\pi_0^{\unsure{\natural}}$ is denoted $\widehat{\unsure{\mathfrak{S}|\mathcal{G}}}(\pi_0^{\unsure{\natural}})$,
it is therefore the inverse image of a subgroup of
$\widehat{\widehat{\mathcal{T}}}(\pi)$, denoted
$\widehat{\mathcal{T}}(\pi_0^{\unsure{\natural}})$. We thus have an inclusion of
extension structures

(10)
\diagram{\begin{tikzcd}[ampersand replacement=\&]
1 \arrow[r] \& \pi \arrow[r] \&
\widehat{\unsure{\mathfrak{S}|\mathcal{G}}}(\pi_0^{\unsure{\natural}}) \arrow[r] \arrow[d,hook] \&
\widehat{\mathcal{T}}(\pi_0^{\unsure{\natural}}) \arrow[r] \arrow[d,hook] \& 1 \\
1 \arrow[r] \& \pi \arrow[r] \&
\widehat{\widehat{\unsure{\mathfrak{S}|\mathcal{G}}}}(\pi) \arrow[r] \&
\widehat{\widehat{\mathcal{T}}}(\pi) \arrow[r] \& 1
\end{tikzcd}}

\margin{\uline{N.B.} The set of classes of discretifications of $\pi$ is described
\illegible{2} of the \uline{outer} group $\pi$, as
$\operatorname{Isom\,ext}(\widehat{\Pi}_{g,\nu},\pi)$\deleted{\illegible{1}}
divided by $\widehat{\mathcal{T}}_{g,\nu}$.
We moreover have \unsure{an arrow} of degree $2$
\diagram{\begin{tikzcd}[ampersand replacement=\&]
\operatorname{Isom\,ext}(\widehat{\Pi}_{g,\nu},\pi)/\widehat{\mathcal{T}}_{g,\nu}^{+}
\arrow[d] \\
\operatorname{Isom\,ext}(\widehat{\Pi}_{g,\nu},\pi)/\widehat{\mathcal{T}}_{g,\nu}
\end{tikzcd}}
which makes it possible for \unsure{every|the} \unsure{class} of discretifications, to define
its two \uline{\unsure{orientations}}, and speak of
\uline{oriented} discretifications.}

\mspage{320}

\scanpage{322}
\mspage{163}

(One shows by arithmetico-geometric means that the inclusion
$\widehat{\mathcal{T}}\longrightarrow\widehat{\widehat{\mathcal{T}}}$
is never an isomorphism.) We thus find a map
\[
\tag{11}
\operatorname{Discr\acute{e}t}^{\unsure{\prime}}(\pi)\longrightarrow
\text{Closed subgroups of }\widehat{\widehat{\mathcal{T}}}(\pi)
\]
obviously compatible with the action of
\deleted{$\unsure{\widehat{\widehat{\mathcal{T}}}}$}
$\widehat{\widehat{\unsure{\mathfrak{S}|\mathcal{G}}}}(\pi)$ (transport of structure) -- operating
on the right via $\widehat{\widehat{\mathcal{T}}}(\pi)$ and its inner automorphisms on itself.

We thus find a conjugacy class, well determined \unsure{of|of the}
closed subgroups of $\widehat{\widehat{\mathcal{T}}}(\pi)$,
which we call its ``\uline{geometric}'' Teichmüller subgroups.
\deleted{It is by no means improbable moreover that
$\widehat{\mathcal{T}}(\pi_0^{\unsure{\natural}})$ is \illegible{1}}
It could moreover be that $\widehat{\mathcal{T}}_{g,\nu}$ is
its own normalizer in $\widehat{\widehat{\mathcal{T}}}_{g,\nu}$
-- equivalent to the assertion that (11) is \unsure{injective|bijective|injection} :
The datum of a class of discretifications\inserted{of $\pi$}
would be equivalent to that of a geometric Teichmüller subgroup in
\deleted{$\unsure{\widehat{\widehat{\mathcal{T}}}}$}
its \unsure{extended} Teichmüller group $\widehat{\widehat{\mathcal{T}}}$.

\margin{this is completely screwy and ultra-false
-- a moment of aberration ! This becomes clear in what follows\ldots
The \unsure{sensible} question (with which \unsure{I} must have
\unsure{at the time} confused) \unsure{is} whether $\widehat{\mathcal{T}}$ is
\uline{invariant} in $\widehat{\widehat{\mathcal{T}}}$, i.e. whether
the normalizer is $\widehat{\widehat{\mathcal{T}}}$ \unsure{all|all} \unsure{of it} !}

\mspage{321}

\scanpage{323}

Note that we have the canonical homomorphisms
\[
\tag{12}
\widehat{\widehat{\mathcal{T}}}\deleted{\unsure{(\pi_0^{\natural})}}
\xrightarrow{\chi}\widehat{\unsure{\mathbb{Z}}}^{*}
\]
\[
\tag{13}
\widehat{\widehat{\mathcal{T}}}\longrightarrow
\unsure{\mathfrak{S}_{I}|\mathcal{G}_{I}}
\]
\margin{where $I=I(\pi)$ is \unsure{the set|the set} of conjugacy classes of loop subgroups of $\pi$}
inducing on $\widehat{\mathcal{T}}(\pi_0^{\unsure{\natural}})$
\deleted{\unsure{continuous}} corresponding homomorphisms -- whence the definition of the
subgroups $\widehat{\widehat{\mathcal{T}}}^{!}$,
$\widehat{\widehat{\mathcal{T}}}^{+}$,
$\widehat{\widehat{\mathcal{T}}}^{!+}$ and likewise for
$\widehat{\mathcal{T}}$. Note that $\chi$ takes on $\widehat{\mathcal{T}}$
only the values $\pm1$.

The subgroup $\widehat{\mathcal{T}}^{+}$ of outer automorphisms of the
profinite \unsure{groups|groups} with loops $\pi$,
\deleted{\illegible{2}}\inserted{fixing} a class of discretifications
$\pi_0^{\unsure{\natural}}$ and with multiplier $+1$, plays a very special role.
In contrast to that \unsure{which one} can conjecture about $\widehat{\mathcal{T}}$
(of which $\widehat{\mathcal{T}}^{+}$ is an \unsure{open|invariant} subgroup of index $2$),
one would rather suppose that $\widehat{\mathcal{T}}^{+}$ is invariant in
$\widehat{\widehat{\mathcal{T}}}$. In any case,
the subgroups thus defined in $\widehat{\widehat{\mathcal{T}}}$
\deleted{which}\inserted{via} \unsure{classes} of discretifications of $\pi$
\deleted{\illegible{1}} (perhaps

\mspage{322}

\scanpage{324}
\mspage{164}

there is only a single and unique one!) will be called the \uline{strict} geometric Teichmüller subgroups. The choice of such a subgroup \inserted{of $\widehat{\widehat{\mathcal{T}}}$} is, in any case, an element of structure \unsure{distinctly} weaker than that of a class of \unsure{discretifications}, and even than that of a geometric Teichmüller subgroup (not strict).

To specify the relations between these two notions, let us first recall that $\widehat{\mathcal{T}}^{+}=\Sigma$ is deduced from $\widehat{\mathcal{T}}$ as the kernel of $\chi\mid\widehat{\mathcal{T}}:\widehat{\mathcal{T}}\longrightarrow\{\pm1\}$. On the other hand, \deleted{let} (for a chosen strict geometric Teichmüller subgroup $\widehat{\mathcal{T}}^{+}$) let us consider
\[
(14)\qquad
\unsure{\mathcal{N}|N}_{\Sigma}=\mathcal{N}
=\operatorname{Norm}_{\widehat{\widehat{\mathcal{T}}}}
\bigl(\deleted{\widehat{\mathcal{T}}^{+}}\inserted{\Sigma}\bigr)
\qquad (\Sigma=\widehat{\mathcal{T}}^{+})
\]
(perhaps always equal to all of $\widehat{\widehat{\mathcal{T}}}$)

obviously (if $\deleted{\widehat{\mathcal{T}}^{+}}\inserted{\Sigma}$ comes from a $\widehat{\mathcal{T}}$)
\[
(15)\qquad
\underset{\inserted{\unsure{=}\widehat{\mathcal{T}}^{+}}}
{\deleted{\widehat{\mathcal{T}}^{+}}\inserted{\Sigma}}
\subset\widehat{\mathcal{T}}\subset\mathcal{N}_{\Sigma}.
\]
\deleted{and $\mathcal{N}/\widehat{\mathcal{T}}^{+}$ is}
The profinite group $\mathcal{N}/\deleted{\widehat{\mathcal{T}}^{+}}\inserted{\Sigma}$
\deleted{is} associated with \deleted{the \unsure{structure}}
$\deleted{\widehat{\mathcal{T}}^{+}}\inserted{\Sigma}$ is denoted
$\boldsymbol{\Gamma}_{\deleted{\unsure{\widehat{\mathcal{T}}^{+}}}\inserted{\Sigma}}$
\deleted{\illegible{1}} \deleted{\illegible{2}}
(if $\deleted{\widehat{\mathcal{T}}^{+}}\inserted{\Sigma}$ is unique, we write $\boldsymbol{\Gamma}_{\pi}$);

\mspage{323}

\scanpage{325}

and the $\widehat{\mathcal{T}}\deleted{\unsure{\mathcal G}}$ \unsure{giving} rise to the same $\deleted{\widehat{\mathcal{T}}^{+}}\inserted{\Sigma}$ correspond \deleted{\unsure{to the classes}} to a conjugacy class of elements of order $2$ of $\boldsymbol{\Gamma}_{\deleted{\unsure{\widehat{\mathcal{T}}^{+}}}\inserted{\Sigma}}$. It is clear \deleted{\unsure{that we|that}} in any case that we have an injective map
\[
(16)\qquad
\begin{gathered}
\text{set of geometric}\\
\text{Teichmüller subgroups}\\
\widehat{\mathcal{T}}\text{ of }\widehat{\widehat{\mathcal{T}}}\text{ such that}\\
\widehat{\mathcal{T}}^{+}=\Sigma
\end{gathered}
\longrightarrow
\begin{gathered}
\text{set of }\unsure{\text{elements}|\text{elements}}\\
\text{of order }2\text{ of }\boldsymbol{\Gamma}_{\Sigma}
\end{gathered}
\quad :
\]
and that its image is stable under conjugation. Let us show that \unsure{two|the} elements of the image are conjugate under $\boldsymbol{\Gamma}_{\Sigma}$. Indeed, let
\[
\widehat{\mathcal{T}},\ \widehat{\mathcal{T}}^{\prime}\supset\Sigma
\quad\text{such that}\quad
\Sigma=\widehat{\mathcal{T}}^{+}=\widehat{\mathcal{T}}^{\prime+}.
\]
\deleted{We \illegible{1} \unsure{obviously}}
\inserted{$\exists\,g\in\widehat{\widehat{\mathcal{T}}}$}
such that
$\widehat{\mathcal{T}}^{\prime}=\operatorname{int}(g)\unsure{\mathcal{T}}$,
and we \unsure{will have|have} therefore
$\widehat{\mathcal{T}}^{\prime+}=\operatorname{int}_{g}\widehat{\mathcal{T}}^{+}$
i.e. $g\in\mathcal{N}$, OK.

The elements of order $2$ thus obtained in $\boldsymbol{\Gamma}_{\Sigma}$ are called the canonical involutions. We will give a remarkable conjectural interpretation of them below. The set $\widehat{\widehat{\mathcal{T}}}/\mathcal{N}$ is identified with the set of \unsure{subgroups} of

\mspage{324}